% Minimal retained preamble for lane-p2 cards taken from
% sources/logic-fl02-conservation-of-ramsey-s-theorem-f/pinned-original-source.tex
% (Conservation of Ramsey's theorem for pairs, arXiv:2402.11616v3).
%
% Copied verbatim from the author's own file:
%   line 2      : \documentclass{article}
%   lines 16-25 : the author's theorem declarations
%   lines 46-88 : the author's macro definitions for the reverse-mathematics principles
%   line 91     : \card
%   lines 104-107: the natural sums and products, which rely on mathabx (source line 103)
%   line 116    : \qt
%
% Omitted author preamble lines and why (nothing was renamed or redefined):
%   line 3 (graphicx), 7 (url), 9 (cancel): no figure, url or cancel use in the retained spans;
%   line 6 (cleveref): the assembler adds its own cleveref machinery, a second load clashes;
%   line 8 (\usepackage[utf8x]{inputenc}): inputenc with the utf8x option fights the forced
%       xelatex engine;
%   lines 28-37 (the rep@theorem machinery and \newreptheorem{maintheorem}): the retained spans
%       use no rep-theorem environment;
%   lines 40-43 (xcolor and the three author comment macros): not used by the retained spans;
%   lines 93-99 (the LS2 font encoding block for \smashtimes based on stix): the host TeX Live has
%       no stix directory and the author himself never uses the result -- line 105 maps \ntimes to
%       \dottimes, and line 110 shows \smashtimes already commented out in the source.
% mathabx is retained because \dotplus, \dottimes and \dotdiv come from it.
\documentclass{article}
\usepackage{amsmath,amssymb,amsthm}
\usepackage{mathabx}
\usepackage{cancel}
\usepackage{cleveref}

\newtheorem{theorem}{Theorem}
\numberwithin{theorem}{section}
\newtheorem{maintheorem}[theorem]{Main Theorem}
\newtheorem{lemma}[theorem]{Lemma}
\newtheorem{question}[theorem]{Question}
\newtheorem{proposition}[theorem]{Proposition}
\newtheorem{statement}[theorem]{Statement}
\newtheorem{remark}[theorem]{Remark}
\newtheorem{definition}[theorem]{Definition}
\newtheorem{corollary}[theorem]{Corollary}

\newcommand{\RCA}{\mathsf{RCA}}
\newcommand{\WKL}{\mathsf{WKL}}
\newcommand{\SRT}{\mathsf{SRT}}
\newcommand{\EM}{\mathsf{EM}}
\newcommand{\SEM}{\mathsf{SEM}}
\newcommand{\ADS}{\mathsf{ADS}}
\newcommand{\SCAC}{\mathsf{SCAC}}
\newcommand{\SADS}{\mathsf{SADS}}
\newcommand{\COH}{\mathsf{COH}}
\newcommand{\ACA}{\mathsf{ACA}}
\newcommand{\qvdash}{\operatorname{{?}{\vdash}}}
\newcommand{\nqvdash}{\operatorname{{?}{\nvdash}}}
\newcommand{\BSig}{\mathsf{B}\Sigma^0}
\newcommand{\ISig}{\mathsf{I}\Sigma^0}
\newcommand{\LSig}{\mathsf{L}\Sigma^0}
\newcommand{\BPi}{\mathsf{B}\Pi^0}
\newcommand{\IPi}{\mathsf{I}\Pi^0}
\newcommand{\LPi}{\mathsf{L}\Pi^0}
\newcommand{\IDel}{\mathsf{I}\Delta^0}
\newcommand{\BCDe}{\mathsf{BC}\Delta^0}
\newcommand{\WF}{\mathsf{WF}}
\newcommand{\RT}{\mathsf{RT}}
\newcommand{\BME}{\mathsf{BME}}
\newcommand{\TT}{\mathsf{TT}}
\renewcommand{\O}{\mathcal{O}}
\renewcommand{\L}{\mathcal{L}}
\newcommand{\C}{\mathcal{C}}
\newcommand{\R}{\mathcal{R}}
\renewcommand{\S}{\mathcal{S}}
\newcommand{\T}{\mathcal{T}}
\renewcommand{\P}{\mathcal{P}}
\newcommand{\M}{\mathcal{M}}
\newcommand{\Nc}{\mathcal{N}}
\renewcommand{\P}{\mathcal{P}}
\newcommand{\Q}{\mathcal{Q}}
\newcommand{\NN}{\mathbb{N}}
\newcommand{\PP}{\mathbb{P}}
\newcommand{\uh}[0]{{\upharpoonright}}
\newcommand{\card}{\operatorname{card}}
\newcommand{\nplus}{\dotplus}
\newcommand{\ntimes}{\dottimes}
\newcommand{\nminus}{\dotdiv}
\newcommand{\nsum}{\dot{\sum}}
\def\qt#1{``#1''}
\title{The set attached to a node of the SEM forcing is transitive}
\author{Quentin Le Houérou, Ludovic Levy Patey, Keita Yokoyama}
\date{arxiv:2402.11616v3; CC BY 4.0}
\begin{document}
\maketitle
\noindent\textit{Editorial note.} The counted unit is the article's own Lemma lem:sem-wf-nodes-transitive at source lines 831-833 together with the proof environment that begins on the next line (834-839); the pairing is the ordinary adjacent one.\par\medskip
\noindent\textit{Editorial note.} The retained context at source line 817 is the author's introduction of the combinatorial concepts used by the two technical lemmas of this section; the retained context at source lines 819-829 is his own definition of a node of a condition, of the union pi of a node, of validity and of satisfaction of a Sigma\textasciicircum{}0\_1 formula.\par\medskip
\noindent\textit{Editorial note.} The retained preamble reproduces verbatim the author's document class at source line 2, his theorem declarations at source lines 16-25, his macro definitions at source lines 46-88 and 91, his natural-operation macros at source lines 104-107 and his quotation macro at source line 116, and it retains mathabx from source line 103 and cleveref, which the assembler only adds by itself when it generates navigation stubs. The LS2 font-encoding block at source lines 93-99 is not retained because the host TeX Live has no stix directory, and the inputenc load with the utf8x option at source line 8 is not retained because it fights the forced xelatex engine.\par\medskip
\noindent\textit{Editorial note.} No retained line contains a citation or a cross-reference, so the material carries no bibliography and needs no navigation pointer.\par\medskip
\noindent\textit{Editorial note.} The article's conservation theorems, its COH and WKL lemmas and the other lemmas of the SEM section are not counted here and are not used as premises.\par\medskip
\noindent\textit{Allowed standard prerequisites (editorial).}\begin{itemize}
\item Conditions of the SEM forcing and the sets X\_a and A\_i(R). In the SEM section the requirements to be satisfied are listed by the author at source lines 810-815, and a condition is a pair (sigma, a) built from a finite set and an infinite set X\_a; the sets A\_i(R) for i below two are the two sides of the partition induced by the stable colouring R on X\_a.. Fixes the objects of the counted lemma, namely a condition (sigma, a) and one of its nodes, and the colouring R with respect to which transitivity is asserted.
\item Node of a condition and the set pi of a node. A (sigma, a)-node is an M-finite sequence of pairs (i\_s, rho\_s) where each i\_s is below two and each rho\_s is a non-empty M-finite R-transitive subset of X\_a whose minimum exceeds the maximum of the previous one, with the comparability pattern between consecutive blocks determined by the bits i\_s; pi of the node is the union of sigma with all the blocks; the author's own definition is retained at source lines 819-829.. The counted lemma asserts that this union is R-transitive, so the object pi and the constraints on the blocks must be available.
\item Minimal interval of a transitive set and Sigma\textasciicircum{}0\_0 induction. For an R-transitive set, the minimal intervals are the maximal blocks of consecutive elements; Sigma\textasciicircum{}0\_0 induction is available in the models considered, which is what permits the author's rank-by-rank argument.. The proof proceeds by Sigma\textasciicircum{}0\_0 induction on the rank of the block and shows that each new block lies inside a minimal interval of the union built so far.
\end{itemize}\medskip
We will need two technical lemmas, which are Shore blocking arguments for satisfying the two kinds of requirements, namely, deciding the jump and preserving $\WF(\alpha)$. Before this, we introduce some basic combinatorial concepts which will be useful for both lemmas. A $\Sigma^0_1$ formula $\varphi(G)$ is in \emph{normal form} if $\varphi(G) \equiv \exists i \psi(G \uh i)$ for a $\Delta^0_0$ formula $\psi$.

\begin{definition}
Fix a condition $(\sigma, a)$.
\begin{enumerate}
    \item[(1)] A \emph{$(\sigma, a)$-node} is an $M$-finite sequence of the form
$$\nu = \langle (i_0, \rho_0), (i_1, \rho_1), \dots (i_{q-1}, \rho_{q-1}) \rangle$$ such that for every~$s < q$, $i_s < 2$, $\rho_s \subseteq X_a$ is a non-empty $M$-finite $R$-transitive set with $\min \rho_s > \max \rho_{s-1}$ (letting $\rho_{-1} = \sigma$), and for every~$s < t < q$, if $i_s = 0$ then $\rho_s \to_R \rho_t$, and if $i_s = 1$ then $\rho_t \to_R \rho_s$.
We then write $$\pi(\nu) = \sigma \cup \rho_0 \cup \dots \cup \rho_{q-1}$$
    \item[(2)] This $(\sigma, a)$-node is \emph{valid} if for every $s < q$, $\rho_s \subseteq A_{i_s}(R)$
%    \item[(3)] A (finite or infinite) set~$Y$ is \emph{compatible} with this $(\sigma, a)$-node if for every if for every~$s < q$, if $i_s = 0$ then $\rho_s \to_R Y$, and if $i_s = 1$ then $Y \to_R \rho_s$.
    \item[(3)] A $(\sigma, a)$-node $\nu$ \emph{satisfies} a $\Sigma^0_1(\M)$-formula $\varphi(G, x) \equiv \exists i \psi(G \uh i, x)$  if $\psi(\pi(\nu), |\nu|-1)$ holds.
\end{enumerate}
\end{definition}

\begin{lemma}\label[lemma]{lem:sem-wf-nodes-transitive}
Fix a $(\sigma, a)$-node $\nu$. Then $\pi(\nu)$ is $R$-transitive.
\end{lemma}

\begin{proof}
Say $\nu = \langle (i_0, \rho_0), (i_1, \rho_1), \dots (i_{q-1}, \rho_{q-1}) \rangle$.
We prove by $\Sigma^0_0$-induction that for every~$\ell < t$, $\rho_\ell$ is an $R$-transitive tournament included in a minimal interval of $\sigma \cup \rho_0 \cup \dots \cup \rho_{\ell-1}$ and that $\sigma \cup \rho_0 \cup \dots \cup \rho_{\ell}$ is $R$-transitive.

This is true for $\ell = 0$ since $(\sigma, a)$ is a condition. And assuming the property to be true at a rank $\ell < t-1$, the construction of $\rho_{\ell +1}$ ensures that it is included in the minimal interval of $(\sigma,a)$ and that $\rho_{\ell + 1} \to_R \rho_{j}$ if $i_j = 1$ and $\rho_j \to_R \rho_{\ell + 1}$ if $i_j = 0$, so no elements of $\sigma \cup \rho_0 \cup \dots \cup \rho_{\ell}$ is between two elements of $\rho_{\ell+1}$. Since $\sigma \cup \rho_0 \cup \dots \cup \rho_{\ell}$ is $R$-transitive (by the induction hypothesis), this means that $\rho_{\ell +1}$ is included in one of its minimal intervals, and since it is itself $R$-transitive, $\sigma \cup \rho_0 \cup \dots \cup \rho_{\ell} \cup \rho_{\ell +1}$ is $R$-transitive.
\end{proof}

\end{document}
